\documentclass{article}

\begin{document}
    \begin{center}
        \Large \textbf{Rabbit Challenge 2: Installiere, konfiguriere und verwende es} \\ [6pt]
        \normalsize
        Autor: Maximilian Jakob Maag B.Sc. \\
        Version: 1.0
    \end{center}

    \section{Intro}
    Playbooks werden häufig dafür verwendet Software zu installieren und zu konfigurieren.

    \section{Aufgaben}
    In dieser Challenge geht es darum, ein Ansible Playbook zu erstellen, dass eine Anwendung installiert.
    Es ist wichtig zu betonen, dass keine Installationsschritte händisch ausgeführt werden dürfen.

    \subsection{Aufgabe 1}
    Deployen Sie eine VM. Legen Sie einen DNS-Eintrag für die VM an, wie er in der Datei README.md beschrieben ist.

    \subsection{Aufgabe 2}
    Stellen Sie sicher das Ihr SSH Public Key hinterlegt ist.

    \subsection{Aufgabe 3}
    Installieren Sie eine PostgreSQL Datenbank auf der VM.
    Richten sie einen Benutzer auf der DB ein. Erstellen Sie die DB nextcloud und
    geben Sie alle Rechte auf die DB an den Benutzer.

    \subsection{Aufgabe 4}
    Installieren Sie PHP und die benötigten Module für Nextcloud.

    \subsection{Aufgabe 5}
    Legen Sie eine Konfigurationsdatei für Apache an.
    Am Port 80 soll später die Nextcloud erreichbar sein.

    \subsection{Aufgabe 6}
    Beziehen Sie ein HTTPS Zertifikat von Let's Encrypt und richten Sie es für die Nextcloud Installation ein.
    Stellen Sie sicher, dass HTTP Anfragen automatisch auf HTTPS umgeleitet werden.

    \subsection{Aufgabe 7}
    Installieren Sie nun Nextcloud selbst. Es muss nativ auf der Maschine installiert werden.
    Das Nutzen von Docker oder ähnlichen Technologien ist nicht erlaubt. Als Webserver kann Apache oder Nginx verwendet werden.

    \subsection{Aufgabe 8}
    Ändern Sie Ihr Playbook so, dass es eine Nextcloud replizierbar installieren kann.
    Führen Sie es anschließend auf mindestens 20 Servern aus. Jeder Server muss einen eigenen DNS Eintrag besitzen.
    Deployen Sie Ihre Server in mindestens 6 verschiedene Regionen. Zum Beispiel: US-east1, US-west1, EU-west1, EU-west2, ASIA-east1, ASIA-east2.

\end{document}